\chapter{Arquitectura de producción}
\label{cap:AnexoI}

Este anexo describe la arquitectura de producción del sistema: un VPS europeo orquestado con Docker Compose y Traefik como proxy inverso. El mismo \emph{stack} que se usa en desarrollo se despliega en producción sin modificaciones, añadiendo únicamente la capa de enrutamiento y TLS que proporciona Traefik.

\section{Motivación y elección de infraestructura}

La elección de un VPS frente a proveedores cloud gestionados responde principalmente al coste. El motor de IA requiere un nodo con GPU para inferencia en tiempo real; en Azure o Google Cloud Platform, un nodo de ese tipo supera los 1.000\,€/mes. Un VPS con CPU moderna en Europa cuesta aproximadamente 130\,€/año, con el motor de IA ejecutándose en modo CPU (latencia inferior a 1\,s por informe, suficiente para el caso de uso).

La simplicidad operativa es el segundo factor. No hay clúster que gestionar, ni \emph{control plane}, ni herramientas adicionales: \texttt{provision.sh} + \texttt{make setup} + \texttt{make deploy} y el sistema está en funcionamiento. El número de piezas que pueden fallar es mínimo.

La portabilidad futura no se compromete: todos los servicios están dockerizados y la única pieza específica del entorno es el fichero \texttt{docker-compose.cpu.yml}. Si en el futuro el volumen de usuarios justifica el coste de un clúster cloud o Kubernetes, la migración no requiere cambios en ningún componente de la aplicación.

\section{Infraestructura}

El servidor es un VPS de Contabo\footnote{\url{https://contabo.com}} ubicado en Europa (Frankfurt). La pila de software se compone de:

\begin{itemize}
    \item \textbf{Sistema operativo}: Ubuntu LTS.
    \item \textbf{Docker Engine}: gestor de contenedores y red.
    \item \textbf{Docker Compose}: orquestación de los servicios de la aplicación.
    \item \textbf{Traefik v3}: proxy inverso, terminación TLS y enrutamiento por dominio.
\end{itemize}

\section{Aprovisionamiento}

El script \texttt{provision.sh} automatiza la puesta en marcha de un servidor desde cero:

\begin{lstlisting}[language=bash, caption={Pasos del script de aprovisionamiento.}, label=lst:provision-sh]
# 1. Crear usuario `app` sin privilegios
# 2. Generar clave SSH ed25519 y añadirla a GitHub (pausa interactiva)
# 3. Instalar Docker Engine via get.docker.com
# 4. Clonar el repositorio en /home/app/CIE-10
# 5. Copiar .env.example -> .env (para edición manual de credenciales)
\end{lstlisting}

Tras ejecutar \texttt{provision.sh}, el operador edita el fichero \texttt{.env} con las credenciales reales (base de datos, SMTP, claves de aplicación, token de HuggingFace) y lanza el despliegue.

\section{Flujo de despliegue}

El proceso completo desde un servidor recién aprovisionado se reduce a tres pasos:

\begin{enumerate}
    \item \textbf{Aprovisionamiento del servidor}: \texttt{sudo bash provision.sh} (requiere root para instalar Docker, crear el usuario \texttt{app} y configurar fail2ban).
    \item \textbf{Setup inicial} (construcción de imágenes, migraciones y seed): \texttt{make setup}
    \item \textbf{Despliegue}: \texttt{make deploy}
\end{enumerate}

\texttt{make deploy} arranca en orden: la red Docker compartida, Traefik (proxy), el \emph{stack} de monitorización (Prometheus + Grafana) y los servicios de la aplicación (base de datos, backend, frontend, motor de IA). En entorno de desarrollo el flujo es equivalente pero sin Traefik ni monitorización: \texttt{make setup up}.

\section{Proxy inverso con Traefik}

Traefik se despliega como un servicio Docker independiente definido en \texttt{docker-compose-proxy.yml}. Su configuración es declarativa: cada servicio de la aplicación expone sus rutas mediante etiquetas Docker (\emph{labels}), y Traefik las descubre automáticamente sin necesidad de recargar configuración.

\begin{lstlisting}[language=bash, caption={Fragmento de configuración de Traefik.}, label=lst:traefik-config]
# Entrypoints: HTTP (:80) redirige permanentemente a HTTPS (:443)
# Let's Encrypt: certificados TLS via HTTP challenge, renovación automática
# Red Docker: escucha en la red `ciecoder` y descubre servicios por labels
# Dashboard: protegido con Basic Auth y rate-limiting
\end{lstlisting}

Los certificados TLS se obtienen y renuevan automáticamente mediante Let's Encrypt (HTTP challenge). El dominio \texttt{clicoder.app} está registrado y apunta al servidor de producción.

\section{Protección con Cloudflare}

El dominio \texttt{clicoder.app} está configurado con el proxy de Cloudflare activado. El tráfico de los usuarios pasa primero por la red de Cloudflare antes de llegar al VPS, lo que aporta tres capas de protección sin coste adicional:

\begin{itemize}
    \item \textbf{Protección frente a bots y \emph{scraping}}: Cloudflare filtra el tráfico automatizado malicioso y bloquea los agentes de rastreo no autorizados antes de que alcancen el servidor.
    \item \textbf{Mitigación de ataques de denegación de servicio} (DDoS): la red distribuida de Cloudflare absorbe el tráfico volumétrico sin que el VPS llegue a recibir la carga.
    \item \textbf{Ocultación de la IP del servidor}: la dirección IP real del VPS no es visible desde el exterior; los atacantes solo conocen las IPs de Cloudflare, lo que reduce la superficie de ataque directa.
\end{itemize}

Esta capa de protección es complementaria a Traefik: Cloudflare actúa en el borde de la red, antes de que el paquete llegue al servidor, mientras que Traefik gestiona el enrutamiento y TLS una vez que el tráfico ya ha pasado por el proxy.

\section{Protección con Fail2ban}

Cloudflare filtra el tráfico en el borde de la red, pero no actúa sobre intentos de acceso por SSH ni sobre peticiones que llegan directamente al servidor. Fail2ban cubre esta superficie residual: monitoriza los logs del sistema y de Traefik, y bloquea a nivel de \texttt{iptables} las IPs que muestran comportamiento de fuerza bruta.

La distinción respecto al \emph{rate limiting} de Traefik es relevante. El \emph{middleware} de Traefik devuelve \texttt{429 Too Many Requests}, pero la IP sigue abriendo conexiones TCP y consumiendo recursos del servidor. Fail2ban, en cambio, inserta una regla en \texttt{iptables} que descarta los paquetes antes de que lleguen a ningún servicio.

El script \texttt{provision.sh} instala y configura fail2ban con dos \emph{jails}:

\begin{table}[h]
\centering
\caption{Jails de fail2ban configurados en producción.}
\label{tab:fail2ban-jails}
\begin{tabular}{llll}
\hline
\textbf{Jail} & \textbf{Fuente de logs} & \textbf{Umbral} & \textbf{Ban} \\
\hline
\texttt{sshd}         & \texttt{/var/log/auth.log}            & 3 intentos fallidos    & 24 h \\
\texttt{traefik-auth} & \texttt{/var/log/traefik/access.log}  & 10 respuestas 401 en 5 min & 1 h  \\
\hline
\end{tabular}
\end{table}

El \emph{jail} de Traefik se apoya en que el servicio escribe los \emph{access logs} en formato JSON al fichero \texttt{/var/log/traefik/access.log}, montado como \emph{bind mount} desde el host. El filtro personalizado \texttt{traefik-auth} extrae el campo \texttt{ClientHost} de cada línea con \texttt{DownstreamStatus: 401}, lo que cubre cualquier \emph{endpoint} de autenticación del backend sin necesidad de hardcodear rutas.

\section{Copias de seguridad}

Contabo incluye en el plan contratado un servicio de \emph{snapshots} diarias automáticas del servidor completo. Cada \emph{snapshot} captura el estado íntegro del sistema, incluyendo los volúmenes Docker donde reside la base de datos PostgreSQL, por lo que no es necesario implementar un mecanismo de backup con \texttt{pg\_dump} adicional.

\section{Comparativa de costes}

La Tabla~\ref{tab:vps-vs-cloud} muestra la diferencia de coste entre el setup actual y dos alternativas cloud gestionadas: una con instancia GPU dedicada y otra íntegramente en CPU.

\begin{table}[h]
\centering
\caption{Comparativa de costes: VPS vs.\ cloud gestionado (GCP, región europea, on-demand, junio 2026).}
\label{tab:vps-vs-cloud}
\begin{tabular}{llr}
\hline
\textbf{Setup} & \textbf{Descripción} & \textbf{Coste aprox.} \\
\hline
VPS (actual)          & VPS CPU, Contabo Frankfurt, Docker Compose    & $\sim$130\,€/año \\
\hline
Cloud con GPU\footnotemark[1]
                      & Nodos CPU (frontend + backend)                & $\sim$140\,€/mes \\
                      & Nodo GPU (motor de IA, 1$\times$ V100)\footnotemark[2] & $\sim$1.950\,€/mes \\
                      & Base de datos gestionada (PostgreSQL)\footnotemark[3] & $\sim$230\,€/mes \\
                      & \textbf{Total estimado}                       & \textbf{$\sim$2.320\,€/mes} \\
\hline
Cloud solo CPU        & Nodos CPU (frontend + backend)                & $\sim$140\,€/mes \\
                      & Nodo CPU para motor de IA (8\,vCPU, 30\,GB)\footnotemark[4] & $\sim$230\,€/mes \\
                      & Base de datos gestionada (PostgreSQL)\footnotemark[3] & $\sim$230\,€/mes \\
                      & \textbf{Total estimado}                       & \textbf{$\sim$600\,€/mes} \\
\hline
\end{tabular}
\end{table}
\footnotetext[1]{GCP europe-west4. Precios on-demand consultados en junio de 2026. Fuente: \url{https://cloud.google.com/compute/gpus-pricing}}
\footnotetext[2]{Instancia n1-standard-8 + 1$\times$ NVIDIA V100: \$2,55/h GPU + \$0,35/h VM $\approx$ \$2.117/mes $\approx$ 1.950\,€/mes.}
\footnotetext[3]{Cloud SQL PostgreSQL, 4\,vCPU / 16\,GB RAM, estimado para europe-west4. Fuente: \url{https://cloud.google.com/sql/pricing}}
\footnotetext[4]{Instancia n1-standard-8 (8\,vCPU, 30\,GB RAM) para alojar el modelo RigoBERTa-Clinical en CPU ($\sim$13\,GB de RAM en inferencia).}

La GPU no es estrictamente necesaria para servir tráfico: el motor de IA ya opera en modo CPU en el VPS actual con latencia inferior a 1\,s por informe. Sin embargo, la opción cloud en CPU sigue siendo mas de 40 veces más cara que el VPS, dado que las instancias de aplicación son pequeñas pero el modelo requiere un nodo con suficiente RAM ($\sim$13\,GB) para cargarse íntegramente. En ambos escenarios cloud el coste está dominado por componentes que en el VPS se comparten en un único nodo sin sobrecoste.
